\documentclass[fontsize=15pt]{jlreq}

%プリアンブル
\usepackage{amsmath}
\usepackage{mathtools}
\pagestyle{empty}
\usepackage[top=15mm, bottom=20mm, left=20mm, right=20mm]{geometry}

\newcommand{\m}[1]{\raisebox{0.08em}{\textcircled{\scriptsize #1}}}



%本文
\begin{document}
\begin{center}
{\LARGE {振り返り}} \\
-生物\m{3}- %単元ごとに変えるか消す
\end{center}
\vspace{1em}
\begin{flushright}
名前：\underline{\hspace{5cm}}
\end{flushright}

% 問題部分
\begin{enumerate}
    \item 生物が長い時間をかけて変化していくことをなんというか．\\
    (\hspace{2cm})
    \vspace{2em} 

    \item 脊椎動物の出現した順に並べよ.また，それぞれの類の例を一つずつ挙げよ．\\
    \begin{itemize}
    \item ( \hspace{2cm}$\rightarrow$ \hspace{2cm} $\rightarrow$\hspace{2cm} $\rightarrow$ \hspace{2cm}) 
    \vspace{2em}
    \item 例( \hspace{2cm}$\rightarrow$ \hspace{2cm} $\rightarrow$\hspace{2cm} $\rightarrow$ \hspace{2cm})
    \end{itemize}
    \vspace{2em} 

    \item 魚類と両生類の幼生は（\hspace{4cm}），両生類の成体,爬虫類，鳥類，哺乳類は（\hspace{4cm}）で呼吸している．両生類は（\hspace{4cm}）でも呼吸している．
    
    \vspace{2em} 

    \item 現在の形やはたらきは違っていても，起源が同じと考えられる器官を何というか．\\
    (\hspace{2cm})
    \vspace{2em} 

    \item 爬虫類と鳥類の中間的な特徴を持つ進化の証拠となる鳥の名前はを何というか．\\
    (\hspace{4cm})
    \vspace{2em} 

\end{enumerate}

\end{document}
